\section{Methodology (Part 2) and Conclusion}

\begin{frame}{Scoring and triage tiers}
\justifying{
\begin{itemize}
  \item Each subgraph is reduced to a fixed feature vector: mean edge
    weight, mean and maximum severity, kill-chain stage coverage, a
    \textit{temporal-order} score, span, and size.
  \item The temporal-order score is the fraction of consecutive
    (time-sorted) events whose MITRE ATT\&CK stage is non-decreasing.
    A genuine attack progresses in order --- Initial Access $\to$
    Privilege Escalation $\to$ Defense Evasion $\to$ Collection $\to$
    Exfiltration --- so it scores $1.0$, whereas uncorrelated noise does not.
  \item An \textbf{Isolation Forest} scores every subgraph without using any
    labels. Its contamination parameter is set from the dataset's
    FP:TP ratio metadata, so the method ports directly to new input files.
  \item Subgraphs are then split into three tiers:
  \begin{itemize}
    \item \textbf{Critical} --- flagged by the model; the analyst must review.
    \item \textbf{Elevated} --- the next $5\%$ by score; review as capacity allows.
    \item \textbf{Low} --- logged, surfaced only on demand.
  \end{itemize}
\end{itemize}}
\end{frame}
\newpage

\begin{frame}{Key performance indicators}
\justifying{
\begin{itemize}
  \item On the provided dataset, the pipeline ingests \textbf{1{,}012} raw
    alerts and decomposes them into \textbf{185} candidate subgraphs.
  \item Only \textbf{8} subgraphs (\textbf{43} events) reach the Critical
    tier --- a \textbf{95.8\%} reduction in the events an analyst must
    triage.
  \item All \textbf{6} ground-truth attack scenarios land in the Critical
    tier: \textbf{100\% recall} at \textbf{75\% precision}. The two
    false positives that accompany them are cheap --- a few minutes of
    review each --- whereas a missed attack is an incident.
  \item Widening to Critical $+$ Elevated still cuts the workload by
    \textbf{90.8\%} while providing a second, lower-priority queue.
\end{itemize}}
\end{frame}
\newpage

\begin{frame}{Performance across datasets}
\justifying{
\begin{itemize}
  \item We ran the identical pipeline on the provided dataset and on two
    independently generated synthetic datasets (same schema, harder
    separation) that carry ground-truth labels.
\end{itemize}
\vspace{0.4em}
\begin{center}
\renewcommand{\arraystretch}{1.3}
\begin{tabular}{lcccc}
\hline
\textbf{Dataset} & \textbf{Events} & \textbf{Subgraphs}
  & \textbf{Recall} & \textbf{Precision} \\
\hline
Provided        & 1{,}012 & 185 & 100\% (6/6) & 75\% \\
Synthetic 1     & 998     & 179 & 83\% (5/6)  & 62\% \\
Synthetic 2     & 997     & 179 & 100\% (6/6) & 75\% \\
\hline
\end{tabular}
\end{center}
\vspace{0.4em}
\begin{itemize}
  \item Recall is measured at the Critical tier. In \textbf{every} dataset,
    \textbf{all 6} attacks fall within Critical $+$ Elevated --- so no
    attack is ever pushed off the analyst's queue.
  \item The single Critical-tier miss (Synthetic 1) is a hard false
    positive that mimics a full kill chain; it still surfaces one tier down.
\end{itemize}}
\end{frame}
\newpage

\begin{frame}{Validating without labels}
\justifying{
\begin{itemize}
  \item Most datasets ship \textbf{no ground-truth labels} --- as in a real
    deployment, where the analyst is the one assigning them. We can still
    sanity-check the output two ways.
  \item \textbf{Against metadata.} The provided unlabeled set declares
    \textbf{3} expected attack scenarios in its metadata. With no access to
    labels, the pipeline flags \textbf{exactly 3} Critical subgraphs ---
    the correct count falls straight out of the FP:TP-driven threshold.
  \item \textbf{Against structure.} The flagged subgraphs share an
    unmistakable profile that noise does not:
\end{itemize}
\vspace{0.3em}
\begin{center}
\renewcommand{\arraystretch}{1.3}
\begin{tabular}{lcc}
\hline
 & \textbf{Critical} & \textbf{Low} \\
\hline
Avg.\ MITRE stages covered & 5.0 / 5 & 2.1 / 5 \\
Avg.\ edge weight          & 1.00    & 0.59 \\
Avg.\ severity             & 7.6     & 2.5 \\
\hline
\end{tabular}
\end{center}
\vspace{0.3em}
\begin{itemize}
  \item Every flagged subgraph spans the \textbf{full kill chain} with
    \textbf{maximal-confidence} edges --- exactly what a true positive
    should look like.
\end{itemize}}
\end{frame}
\newpage

\begin{frame}{Making it legible: subgraph visualization}
\justifying{
\begin{itemize}
  \item \texttt{visualize\_graphs.py} renders each subgraph as a
    \textbf{bipartite graph}: event nodes (coloured by MITRE stage) on one
    side, the users, hosts, files, and domains they touch on the other.
  \item Edge thickness encodes the relationship \textbf{weight}, so a
    high-confidence attack subgraph is visibly denser and more strongly
    connected than a noise cluster --- the analyst sees \textit{why} it
    scored highly at a glance, with no need to read the raw JSON.
  \item The same tool plots the kill-chain timeline and the
    feature distributions, giving a fast, visual second opinion alongside
    the numeric score for any subgraph the analyst wants to inspect.
\end{itemize}}
\end{frame}
\newpage

\begin{frame}{Ranking, not just flagging}
\justifying{
\begin{itemize}
  \item The tier threshold is a convenience, not a gate. Every subgraph
    receives a continuous anomaly score, and the system prints the
    \textbf{top-10 most anomalous} subgraphs regardless of tier.
  \item This matters because the threshold is deliberately conservative.
    An analyst scanning the ranked list can always look one or two places
    \textit{past} the cut-off and investigate anything that stands out to
    them --- a subgraph does not have to be formally flagged to be reviewed.
  \item Nothing is ever hidden: ``Low'' simply means low-priority, not
    discarded. The full ranked list, with the score explanation for each
    subgraph, remains available.
  \item In practice the genuine attacks occupy the very top of the ranking
    (ranks 1--6 on the provided data), so a human reading top-down reaches
    every real incident before any noise.
\end{itemize}}
\end{frame}
\newpage

\begin{frame}{Manpower and cost saved (per batch)}
\justifying{
\begin{itemize}
  \item Industry estimates put average alert triage at roughly
    \textbf{7 minutes} per alert and a fully-loaded SOC analyst at about
    \textbf{\$67/hour} ($\sim$\$48/hr U.S.\ base $\times\,1.4$ for
    benefits and overhead).\footnotemark
  \item Manually triaging all \textbf{1{,}012} alerts $\approx$
    \textbf{118 analyst-hours} $\approx$ \textbf{\$7{,}900} per weekly batch.
  \item Reviewing only the \textbf{43} Critical-tier events takes
    \textbf{$\sim$5 hours} ($\sim$\$340) --- a saving of
    \textbf{$\sim$113 hours} and \textbf{$\sim$\$7{,}600} per week.
\end{itemize}
\vspace{0.4em}
\begin{center}
\renewcommand{\arraystretch}{1.3}
\begin{tabular}{lccc}
\hline
\textbf{Approach} & \textbf{Alerts} & \textbf{Hours} & \textbf{Loaded cost} \\
\hline
Manual (everything)   & 1{,}012 & 118 & \$7{,}900 \\
Critical only         & 43      & 5   & \$340 \\
Critical $+$ Elevated & 93      & 11  & \$730 \\
\hline
\end{tabular}
\end{center}
\footnotetext{Triage time: Prophet Security, \textit{SOC Capacity Modeling}
(2025). Salary: Glassdoor / Salary.com (2026).}}
\end{frame}
\newpage

\begin{frame}{Manpower and cost saved (annualized)}
\justifying{
\begin{itemize}
  \item The provided dataset covers roughly \textbf{one week} of alerts.
    Extrapolating the same volume across a year ($\times 52$):
\end{itemize}
\vspace{0.4em}
\begin{center}
\renewcommand{\arraystretch}{1.3}
\begin{tabular}{lcc}
\hline
\textbf{Approach} & \textbf{Analyst-hours/yr} & \textbf{Loaded cost/yr} \\
\hline
Manual (everything) & $\sim$6{,}100 & $\sim$\$413{,}000 \\
Critical only       & $\sim$260     & $\sim$\$18{,}000 \\
\hline
\textbf{Saved}      & $\sim$5{,}900 & $\sim$\$395{,}000 \\
\hline
\end{tabular}
\end{center}
\vspace{0.4em}
\begin{itemize}
  \item That is roughly \textbf{3 full-time analysts} freed from manual
    triage --- redeployable to deeper investigation and threat hunting ---
    while \textbf{every known attack is still caught}.
  \item Figures are a first-order extrapolation from one week and assume a
    constant alert rate; they are intended to convey scale, not a budgeted
    forecast.
\end{itemize}}
\end{frame}
\newpage

\begin{frame}{Future work: analyst feedback and retraining}
\justifying{
\begin{itemize}
  \item The current model is fully \textbf{unsupervised} --- it needs no
    labels, which is what lets it run on any new input file on day one.
  \item A natural extension is a \textbf{feedback loop}. Each subgraph in the
    ranked output could expose a one-click tag --- \texttt{TRUE\_POSITIVE},
    \texttt{FALSE\_POSITIVE}, or \texttt{NEEDS\_INVESTIGATION} --- written to
    a standardized log alongside the subgraph's feature vector.
  \item Because we already compute a feature vector per subgraph, every tag
    becomes a labelled training example essentially for free. Once enough
    accumulate, they support a \textbf{supervised} layer (e.g.\ logistic
    regression or a random forest) that complements the anomaly score and
    learns this organisation's specific notion of a true threat.
  \item Tags would also let the threshold \textbf{self-calibrate}: repeated
    confirmations just below the cut-off can nudge the contamination
    estimate automatically. This was scoped but left as future work.
\end{itemize}}
\end{frame}
\newpage

\begin{frame}{Conclusion}
\justifying{
\begin{itemize}
  \item We turn a flat stream of \textbf{1{,}012} alerts into a ranked,
    tiered queue, cutting the analyst's triage load by roughly \textbf{96\%}
    --- on the order of \textbf{\$0.4M} and \textbf{3 analyst-FTEs} per year
    --- while catching \textbf{every} known attack on the provided data.
  \item MITRE ATT\&CK kill-chain \textbf{temporal logic} is the key signal:
    it distinguishes coordinated attacks from coincidental clusters of
    high-severity noise, and it is what carries performance over to harder,
    unseen data.
  \item The system is unsupervised at its core and ranking-based at the
    surface: nothing is hidden, the full ordered list is always available,
    and analyst feedback offers a clear path to a self-improving model.
\end{itemize}}
\end{frame}